\section{A complete norm from reciprocal-integer samples (NS-78)}
\label{sec:ns78-sampling}

\paragraph{Purpose.}
NS-75 identifies the signed divisor defects but shows that cancelling
individual jumps need not decrease the norm. Here we retain the entire
residual through its reciprocal-integer values, with constants independent
of the number and size of the coefficients. This is an exact geometric
reduction and a complete physical-space replay, not the missing arithmetic
estimate. The infinite sample sum is not replaced by a finite prefix.

For arbitrary real $\alpha,c_1,\ldots,c_N$, put
\[
 R=\alpha\tau-\sum_{n\leq N}c_n\sigma_n,\qquad
 \eta=\sum_{n\leq N}\frac{c_n}{n},\qquad
 z_m=R(1/m),\qquad w_m=\frac1{m(m+1)}.
\]
Define the complete nonnegative arithmetic quantity
\begin{equation}
 \mathcal S(c)=\eta^2+
       \sum_{m=1}^\infty w_m(z_m^2+z_{m+1}^2).
 \label{eq:ns78-sampling-functional}
\end{equation}
Its samples are explicitly
\[
 z_m=-\eta m+\sum_{k\leq m}
 \left(\alpha\mathbf1_{k=1}+\sum_{\substack{n\mid k\\n\leq N}}c_n\right)
                    \log(m/k).
\]
All cancellation inside each displayed sample is retained before squaring.

\begin{proposition}[Uniform complete sampling comparison]
\label{prop:ns78-sampling}
For every real $\alpha$ and finite real coefficient vector,
\begin{equation}
 \frac1{16}\mathcal S(c)\leq\|R\|_2^2
                 \leq\frac{21}{16}\mathcal S(c).
 \label{eq:ns78-sampling-comparison}
\end{equation}
Both sides include the exterior tail and every reciprocal cell. The
infinite series in~\eqref{eq:ns78-sampling-functional} converges.
\end{proposition}
\begin{proof}
Change variables to $t=1/x$ and write $r(t)=R(1/t)$, with measure
$dt/t^2$. On $0<t<1$, $r(t)=-\eta t$, so the complete exterior
energy is exactly $\eta^2$. By NS-75, on $m\leq t\leq m+1$,
\[
 r(t)=-\eta t+D_m\log t-L_m,
\quad D_m=\sum_{k\leq m}\delta_k,
\quad L_m=\sum_{k\leq m}\delta_k\log k.
\]
Here $\delta_k=\alpha\mathbf1_{k=1}+\sum_{n\mid k,\,n\leq N}c_n$
is the signed divisor defect.
Let $h_m=\log(1+1/m)$, $\theta=\log(t/m)/h_m$, and set
\[
 p(t)=(1-\theta)z_m+\theta z_{m+1},\qquad
 b(t)=\theta-(t-m).
\]
Exactly $r(t)=p(t)+\eta b(t)$ on every cell.
Under $t=m e^{h_m\theta}$ the measure is
$(h_m/m)e^{-h_m\theta}d\theta$. Its density divided by $w_m$
is between $m h_m$ and $(m+1)h_m$, hence between $1/2$ and $2$.
Also
\[
 \int_0^1((1-\theta)a+\theta d)^2d\theta
       =\frac{a^2+ad+d^2}{3}
       \in\left[\frac{a^2+d^2}{6},\frac{a^2+d^2}{2}\right].
\]
Thus, writing $V=\int_1^\infty p(t)^2dt/t^2$ and
$S=\sum_{m\geq1}w_m(z_m^2+z_{m+1}^2)$, we get
$S/12\leq V\leq S$, initially for finite unions of cells.

Concavity of $\log(m+v)$ for $0\leq v\leq1$, and its second
derivative bound, give
\[
 0\leq b(m+v)\leq\frac{v(1-v)}{2m^2h_m}
       \leq\frac1{8m^2h_m}\leq\frac1{4m}.
\]
Consequently
\[
 \Gamma:=\int_1^\infty b(t)^2\frac{dt}{t^2}
 \leq\frac1{16}\sum_{m\geq1}\frac{w_m}{m^2}
 \leq\frac1{16},
\]
since $\sum w_m=1$. The original finite combination lies in $L^2$,
so $p=r-\eta b$ is square integrable. Monotone passage through the
finite-cell estimates now proves finiteness of $S$ and their complete
versions; no convergence of $S$ was assumed.

Cauchy--Schwarz and Young's inequality give
\[
 |2\eta\langle p,b\rangle|
       \leq\eta^2/4+4\Gamma V\leq(\eta^2+V)/4.
\]
Using the exact identity
$\|R\|_2^2=\eta^2+V+2\eta\langle p,b\rangle+\eta^2\Gamma$,
we obtain
\[
 \frac34(\eta^2+V)\leq\|R\|_2^2
 \leq\frac{21}{16}\eta^2+\frac54V
 \leq\frac{21}{16}(\eta^2+V).
\]
Inserting $S/12\leq V\leq S$ proves the proposition.
\end{proof}

\paragraph{The arithmetic input in explicit coordinates.}
Specialize to $\alpha=1$. By the fixed-order smoothed Nyman--Beurling
criterion in NS-61, RH is
equivalent to the existence of coefficient vectors, of unbounded allowed
length, with $\mathcal S(c)\to0$. For the actual optimized coefficients,
this is equivalent to $E_N\to0$. Boundedness of $\mathcal S(c)$ alone
does not suffice; it is not the bounded-Weil-floor criterion.
The local geometric comparison is unconditional, but no sequence with
vanishing complete $\mathcal S$ has been produced. A finite collection
of small samples cannot replace its infinite sum or its cofinal limit.

The case $\alpha=0$ also bounds the complete energy of any finite
correction vector. Its comparison quantity still contains all signed
arithmetic samples; this is not an all-vector comparison with the
smooth repaired model excluded in NS-70.

\paragraph{Exact inversion and the remaining arithmetic constraint.}
The samples also recover the coefficients without a Gram solve. Put
$D_0=0$. Continuity at every reciprocal knot gives
\[
 \eta=-z_1,\qquad
 D_m=\frac{z_{m+1}-z_m+\eta}{\log(1+1/m)},\qquad
 \delta_m=D_m-D_{m-1}.
\]
M\"obius inversion then gives, extending $c_n=0$ beyond $N$,
\begin{equation}
 c_n=\sum_{d\mid n}\mu(n/d)\delta_d-\alpha\mu(n).
 \label{eq:ns78-sample-inverse}
\end{equation}
For $n>N$ the right-hand side must therefore vanish. These infinitely
many arithmetic compatibility conditions distinguish admissible sample
sequences from arbitrary small knot values.

In particular, norm convergence of residuals with $\alpha=1$ forces
$c_n\to-\mu(n)$ at each fixed index $n$. Indeed
$|z_m|\leq4\sqrt{m(m+1)}\,\|R\|_2$ follows from the complete
sampling comparison, while $|\eta|\leq\|R\|_2$ follows from the
exterior tail. The finite formulas above then force $D_m,\delta_m\to0$
for every fixed $m$, and~\eqref{eq:ns78-sample-inverse} gives the claim.
This necessary coefficientwise limit supplies neither a uniform growing-index
bound nor a norm-convergent sharp M\"obius truncation; the latter was
excluded in NS-62/66. No interchange of those two limits is made.

\subsection{Complete cell integrals and explicit tails}
\label{subsec:ns78-tails}
For $m\geq1$, put $a=\log m$, $d=\log(m+1)$ and
\[
 I_{1,m}=\frac{a+1}{m}-\frac{d+1}{m+1},\qquad
 I_{2,m}=\frac{a^2+2a+2}{m}-\frac{d^2+2d+2}{m+1}.
\]
The exact cell contribution is
\begin{equation}
 \int_m^{m+1}\frac{r(t)^2}{t^2}dt
 =\eta^2+D_m^2I_{2,m}+L_m^2w_m
 -\eta D_m(d^2-a^2)+2\eta L_m(d-a)-2D_mL_mI_{1,m}.
 \label{eq:ns78-exact-cell}
\end{equation}
The signed cross terms are evaluated before accumulating cell energies.

We also need a complete tail, since $m$ extends indefinitely even when
$c$ has finite support. Write
\[
 u=\alpha-\tfrac12\sum c_n,\qquad
 v=\tfrac12\sum c_n\log n-\tfrac12\log(2\pi)\sum c_n,
 \qquad D=\tfrac16\sum n|c_n|.
\]
For $t\geq N$,
\begin{equation}
 r(t)=u\log t+v+e(t),\qquad |e(t)|\leq D/t.
 \label{eq:ns78-complete-asymptotic}
\end{equation}
Indeed $A'(z)=\{z\}/z$ almost everywhere and Stirling's formula at
integer endpoints identifies the limiting constant
$A(z)-\tfrac12\log z\to\tfrac12\log(2\pi)$.
Using the periodic Bernoulli polynomials,
\[
 A(z)-\tfrac12\log(2\pi z)
 =-\int_z^\infty\frac{B_1(\{t\})}{t}dt
 =\frac{B_2(\{z\})}{2z}
       -\frac12\int_z^\infty\frac{B_2(\{t\})}{t^2}dt.
\]
The first integral is an improper convergent integral; the second is
absolutely convergent. Since $|B_2|\leq1/6$, the remainder is at most
$1/(6z)$. The periodic $B_2$ is continuous at integer joins, so no
endpoint mass is discarded. Summing this bound gives
\eqref{eq:ns78-complete-asymptotic}.

Fix an integer $M\geq N$ and put $q=u\log M+v$. The complete
main-tail energy and remainder budget are
\[
 T_M=\frac{q^2+2uq+2u^2}{M},\qquad
 U_M=\frac{D^2}{3M^3}.
\]
Thus the actual integral from $M$ to infinity differs from $T_M$ by
at most $2\sqrt{T_MU_M}+U_M$. For the infinite sample tail, set
$Q=|q|+(D+|u|)/M$. Then
\begin{equation}
 \sum_{m\geq M}w_m(z_m^2+z_{m+1}^2)
 \leq \frac{2(Q^2+2|u|Q+2u^2)}{M}.
 \label{eq:ns78-sampling-tail}
\end{equation}
To verify this last bound, on each $m\leq t\leq m+1$ both endpoint
values have absolute value at most $Q+|u|\log(t/M)$:
use~\eqref{eq:ns78-complete-asymptotic} and
$\log((m+1)/t)\leq1/M$. Integrating the square with measure $dt/t^2$
and summing all cells proves~\eqref{eq:ns78-sampling-tail}.
These are bounds on the complete tails, not estimates based on a zero
prefix or an unmeasured remainder.

\paragraph{Independent complete physical-space check.}
For the actual optimized coefficients at $N=16,64,256$, the interval
calculation accumulates every cell through $M=65536$ and encloses the
entire remaining physical and sample tails by the formulas above.
At both 256 and 384 bits, the resulting complete physical norm encloses
the independently evaluated full Gram norm. The complete ratios satisfy
\[
 \mathcal S(c_{16})/E_{16}\in(3.323,3.325),\qquad
 \mathcal S(c_{64})/E_{64}\in(2.9002,2.9004),\qquad
 \mathcal S(c_{256})/E_{256}\in(2.644,2.646).
\]
The analytic physical-tail uncertainty divided by $E_N$ is below
$1.64\cdot10^{-6}$, $5.20\cdot10^{-6}$ and $2.38\cdot10^{-4}$,
respectively. Replaying $N=256$ with $M=131072$ and the doubled complete
Gram-kernel cutoff reduces that tail ratio below $4.17\cdot10^{-5}$.
All solves retain Arb balls. The precise finite part, main tail and
remainder radius are archived separately; the printed combined decimal
balls may be wider due to decimal serialization. Thirty exact Maxima
checks cover the cell, mass, interpolation and tail identities.
This is an independent norm evaluation on finite coefficient vectors,
not a cofinal estimate of their sample sums.

\section{Balancing the exterior tail of divisor feedback (NS-79)}
\label{sec:ns79-balanced-feedback}

The sampling norm isolates the exterior coefficient $\eta$. We therefore
test whether removing its change repairs the failed NS-75 step. This is
an explicitly modified correction, not a reinterpretation of that test.
Let $a_n^{(p)}=-r_N(n)\log^p(2N/n)$ be the three original new-block
physical directions, with the constant convention for $p=0$. Define
\begin{equation}
 \widetilde a_n^{(p)}=a_n^{(p)}\quad(N<n<2N),\qquad
 \widetilde a_{2N}^{(p)}=-2N\sum_{n=N+1}^{2N-1}\frac{a_n^{(p)}}n.
 \label{eq:ns79-balance}
\end{equation}

\begin{proposition}[Exact tail balance and its scope]
\label{prop:ns79-balance}
Each raw correction $f_p=\sum_{n=N+1}^{2N}\widetilde a_n^{(p)}\sigma_n$
vanishes for $x>1/(N+1)$. Thus the held-old step preserves $\eta$.
For $p=0$ it cancels the new derivative jumps for $N<m<2N$, leaving
the jump at $m=2N$ unconstrained. It is not an exact cancellation of
the entire next block.

The original three-direction span augmented by the physical endpoint
atom $\sigma_{2N}$ equals the balanced three-direction span augmented
by that atom. This identity persists after the complete old-space
projection. No assertion of norm descent follows from these identities.
\end{proposition}
\begin{proof}
The weighted coefficient sum in~\eqref{eq:ns79-balance} is zero.
For $x>1/(N+1)$ every new atom equals $1/(nx)$, proving vanishing
and tail preservation. In the open block each proper divisor is old,
so the NS-75 jump argument applies at every index except the modified
last one. More explicitly, for $N<t<2N$ the held-old updated residual
has one formula $-\eta t+D_N\log t-L_N$, with $D_N,L_N$ from NS-78;
no new jumps occur inside this open interval.

The difference between each original and balanced physical vector is
a multiple of the same endpoint coordinate. Hence augmenting either
span by that coordinate gives the same space; applying any linear
old-space projection preserves this equality. In adjacent-difference
coordinates the endpoint atom has vector $v_n=n/(2N)$, modulo its
old boundary term. The transformation is retained exactly in the
finite calculation.
\end{proof}

Old-space reoptimization can change $\eta$ again. The projected gains
below therefore do not preserve the raw tail balance. Conversely,
repeating only raw tail-preserving steps with a fixed nonzero initial
$\eta$ can never reduce squared error below $\eta^2$, because the
exterior tail remains unchanged. This observation does not exclude
initialization with $\eta=0$ or algorithms that reoptimize old coefficients.

\paragraph{Certified finite outcome.}
At 256/384 bits, with a doubled complete-kernel cutoff at $N=256$,
the fractions of full projected gain have the following strict enclosures.
\begin{center}
\begin{tabular}{rccc}
\hline
$N$&Balanced three-span&Endpoint-augmented four-span&Diagonal curvature\\
\hline
16&(.5960,.5961)&(.8542,.8543)&(.9308,.9310)\\
32&(.6580,.6581)&(.8240,.8241)&(.9040,.9041)\\
64&(.8978,.8979)&(.9192,.9193)&(.9783,.9784)\\
128&(.8809,.8810)&(.8810,.8811)&(.9852,.9854)\\
256&(.5996,.5997)&(.7118,.7119)&(.9763,.9764)\\
\hline
\end{tabular}
\end{center}
Thus removing the tail change improves the $N=256$ three-span fraction
from $56.02\%$ to $59.97\%$, and the full four-span reaches $71.18\%$;
both remain below curvature's $97.63\%$. The raw untapered unit step
at $N=256$ still increases squared error by a factor in $(53.07,53.09)$,
although this is smaller than NS-75's factor above $1495$.
Thirteen of the fifteen raw balanced unit steps increase error. The two
exceptions are the linear and quadratic tapers at $N=128$, whose error
ratios lie in $(.9643,.9645)$ and $(.9890,.9893)$. This is neither a
uniform-descent theorem nor an exclusion of all balanced feedback.
Sixteen exact Maxima checks cover endpoint compensation, coordinate
conversion and the complete exterior floor. The four-dimensional span
keeps all cross terms; no mixed term is discarded to claim positivity.

\section{A uniform preconditioner that retains arithmetic (NS-80)}
\label{sec:ns80-arithmetic-preconditioner}

The complete sampling comparison also bounds every finite Gram matrix.
Its comparator retains the arithmetic values $A(m/n)$ at \emph{all}
integer knots; it is not the smooth model from NS-70. This gives a
uniform finite-block efficiency theorem, but neither a fast inverse
nor the cofinal lower bound on available gain needed for RH.

Define the complete arithmetic sampling Gram
\begin{equation}
 Q_{ij}=\frac1{ij}+\sum_{m=1}^\infty\frac{
 A(m/i)A(m/j)+A((m+1)/i)A((m+1)/j)}{m(m+1)}.
 \label{eq:ns80-sampling-gram}
\end{equation}
The series converges by NS-78 and Cauchy--Schwarz. For any finite
coefficient vector, its quadratic form is exactly the complete
sampling quantity for the corresponding pure atom combination.
Thus on every finite index set,
\begin{equation}
 \frac1{16}Q\preceq G\preceq\frac{21}{16}Q.
 \label{eq:ns80-full-comparison}
\end{equation}
Both matrices are positive definite: zero physical norm forces zero
samples and exterior coefficient, and the exact inversion in NS-78
then forces every coefficient to vanish.

\begin{proposition}[Uniform efficiency with the full arithmetic comparator]
\label{prop:ns80-uniform-efficiency}
Use the ordered basis consisting of the old $N$ atoms and the new
adjacent differences. Let $H_G,H_Q$ be the respective Schur complements
of the old blocks in their complete Grams. Then
\[
 \frac1{16}H_Q\preceq H_G\preceq\frac{21}{16}H_Q.
\]
For any nonzero \emph{actual} residual correlation vector $g$, set
$v=H_Q^{-1}g$. Its optimally scaled true gain satisfies
\begin{equation}
 \frac{(g^Tv)^2}{v^TH_Gv}
 \geq\frac{21}{121}\,g^TH_G^{-1}g.
 \label{eq:ns80-gain-fraction}
\end{equation}
This is uniform in $N$ and $g$. All old projections and all cross terms
are included. It is a fraction of the full available block gain, not
of the current residual energy.
\end{proposition}
\begin{proof}
The full comparison survives the invertible change to adjacent
coordinates. For each new vector $v$, the Schur quadratic form is the
minimum of the complete quadratic form over old coordinates. Taking
those minima on both sides preserves the same two constants; the old
minimizers need not agree.

More generally put $a=1/16$, $b=21/16$ and
$M=H_Q^{-1/2}H_GH_Q^{-1/2}$. Its eigenvalues lie in $[a,b]$.
With $s=H_Q^{-1/2}g$, normalize its squared spectral coordinates to
probability weights and write $x=\mathbb E\lambda$,
$y=\mathbb E\lambda^{-1}$. For $a\leq\lambda\leq b$,
\[
 \lambda+\frac{ab}{\lambda}\leq a+b,
 \qquad x+aby\leq a+b,
 \qquad xy\leq\frac{(a+b)^2}{4ab}.
\]
The final inequality follows from $(x-aby)^2\geq0$.
The ratio of the selected gain to the full gain is precisely
$1/(xy)$, hence is at least $4ab/(a+b)^2=21/121$.
This is the usual elementary Kantorovich argument, displayed here
to make every hypothesis and normalization explicit.
\end{proof}

\paragraph{Repeated steps within a fixed finite block.}
Start at $x_0=0$ and minimize the exact block quadratic
$\Phi(x)=x^TH_Gx-2g^Tx$. At each step use the current residual
$r_k=g-H_Gx_k$, direction $p_k=H_Q^{-1}r_k$, and the exact line search
\[
 x_{k+1}=x_k+\frac{r_k^Tp_k}{p_k^TH_Gp_k}p_k,
\]
with termination if $r_k=0$. The remaining block gain is
$r_k^TH_G^{-1}r_k$, so Proposition~\ref{prop:ns80-uniform-efficiency}
reduces it by at least $21/121$ each time. After $k$ steps the captured
fraction of full gain is therefore at least
\[
 1-(100/121)^k.
\]
In particular twenty ideal steps capture more than $97.79\%$, uniformly
in block size. This statement assumes application of the complete
$H_Q^{-1}$ and the true $H_G$; no implementation or complexity bound for
that infinite-sample comparator is asserted. Truncating its sum without
a certified tail would invalidate the theorem's premises.

For completeness, minimizing the full sampling objective over the first
$N$ coefficients gives a residual with squared error at most $21E_N$:
apply both bounds of NS-78 before and after the two minimizations.
This uniform approximation factor is also not a decay estimate for $E_N$.

\paragraph{The remaining obstruction.}
Even perfect capture of every block gain does not prove that the limiting
error is zero. The full block gains telescope to $E_N-E_\infty$, and
NS-76 exhibits a component invisible to all of them. A lower estimate
for available gain relative to the entire current error remains essential.
The present theorem preserves arithmetic in its comparator and supplies a
uniform finite-block efficiency guarantee; it does not provide that lower
estimate, a fast computable replacement for $H_Q$, a bound for the earlier
smooth curvature rule, G2 or RH.
